%--------------------------------------------------------------------%
% LAMPIRAN 1--16
% Untuk menempelkan dokumen PDF, gunakan \includepdf[pages=-]{lampiran/nama-berkas.pdf}
% (paket pdfpages sudah dimuat). Untuk gambar: \includegraphics.
%--------------------------------------------------------------------%

\judulhalaman{LAMPIRAN 1. DOKUMEN LEGALITAS MITRA}
\petunjuk{Akte Pendirian, SK Lembaga, dan NPWP Badan Usaha mitra (Pedoman TA butir 3.3.3). \tahap{\wajib{WAJIB}}{\wajib{WAJIB}}}
\isian{Tempelkan dokumen pada halaman ini}

\judulhalaman{LAMPIRAN 2. SURAT PERNYATAAN KESEDIAAN MITRA}
\petunjuk{Ditandatangani \textit{Top Management} mitra, sebelum seminar proposal (Pedoman TA butir 3.3.3). \tahap{\wajib{WAJIB}}{\wajib{WAJIB}}}
\isian{Tempelkan dokumen pada halaman ini}

\judulhalaman{LAMPIRAN 3. SURAT PERNYATAAN KESANGGUPAN PEMBAGIAN KERJA TIM DAN TIMELINE}
\petunjuk{Formulir Lampiran 6 pedoman, memuat judul Laporan TA individu tiap anggota, ditandatangani seluruh anggota kelompok (Pedoman TA butir 3.3.2 dan 5.1.3). \tahap{\wajib{WAJIB}}{\wajib{WAJIB}}}
\isian{Tempelkan dokumen pada halaman ini}

\judulhalaman{LAMPIRAN 4. EVALUASI RESMI MITRA ATAU BERITA ACARA SERAH TERIMA}
\petunjuk{Dokumen evaluasi atau uji penerimaan yang ditandatangani \textit{Top Management} mitra. \tahap{\tidakperlu{Tidak diperlukan}}{\wajib{WAJIB}}}
\isian{Tempelkan dokumen pada halaman ini}

\judulhalaman{LAMPIRAN 5. FORM PERNYATAAN SELESAI TA}
\petunjuk{Hanya bila mitra tidak dapat hadir pada sidang; diverifikasi pembimbing dan Koordinator TA serta dicatat dalam berita acara (Pedoman TA butir 3.3.3). \tahap{\tidakperlu{Tidak diperlukan}}{\wajib{WAJIB bila berlaku}}}
\isian{Tempelkan dokumen pada halaman ini}

\judulhalaman{LAMPIRAN 6. LOGBOOK INDIVIDU}
\petunjuk{Satu \textit{logbook} per anggota, konsisten dengan riwayat versi dan artefak (Pedoman TA butir 3.3.2). \tahap{\opsional{Opsional}}{\wajib{WAJIB}}}
\isian{Tempelkan dokumen pada halaman ini}

\judulhalaman{LAMPIRAN 7. PEER ASSESSMENT DAN EVALUASI KONTRIBUSI OLEH PEMBIMBING}
\petunjuk{\textit{Peer assessment} diisi tiap anggota; evaluasi pembimbing diisi terpisah (Pedoman TA butir 3.3.4). \tahap{\tidakperlu{Tidak diperlukan}}{\wajib{WAJIB}}}
\isian{Tempelkan dokumen pada halaman ini}

\judulhalaman{LAMPIRAN 8. BUKTI PENGUJIAN DAN VERIFIKASI TKT}
\petunjuk{Berkas hasil uji, bukti TKT minimal 3, dan validasinya. \tahap{\tidakperlu{Tidak diperlukan}}{\wajib{WAJIB}}}
\isian{Tempelkan dokumen pada halaman ini}

\judulhalaman{LAMPIRAN 9. KODE SUMBER, PETA KODE, DAN DOKUMENTASI TEKNIS}
\petunjuk{Tautan repositori dan versinya, dokumentasi untuk demonstrasi, serta peta kode untuk \textit{code walkthrough}: modul atau berkas utama, fungsinya, pembuatnya, dan keputusan teknis beserta alasannya. \tahap{\opsional{Opsional}}{\wajib{WAJIB}}}
\isian{Tempelkan dokumen pada halaman ini}

\judulhalaman{LAMPIRAN 10. LINIMASA PROYEK}
\petunjuk{Rencana dan realisasi tiap tahap, disetujui pembimbing. \tahap{\wajib{WAJIB (kolom rencana)}}{\wajib{WAJIB (lengkap)}}}
\begingroup\footnotesize\singlespacing\renewcommand{\arraystretch}{1.6}
\begin{center}
\begin{tabular}{|>{\raggedright\arraybackslash}p{0.3\textwidth}|p{0.13\textwidth}|p{0.13\textwidth}|p{0.12\textwidth}|p{0.1\textwidth}|}
	\hline
	\textbf{Tahap} & \textbf{Rencana} & \textbf{Realisasi} & \textbf{Bukti} & \textbf{Paraf}\\ \hline
	Seminar proposal & & & & \\ \hline
	Analisis dan perancangan & & & & \\ \hline
	Implementasi prototipe & & & & \\ \hline
	Pengujian dan evaluasi & & & & \\ \hline
	Evaluasi dan serah terima mitra & & & & \\ \hline
	Sidang akhir & & & & \\ \hline
\end{tabular}
\end{center}\endgroup

\judulhalaman{LAMPIRAN 11. BERITA ACARA SEMINAR PROPOSAL}
\petunjuk{Berita acara seminar proposal beserta lembar masukan. \tahap{\wajib{WAJIB}}{\tidakperlu{Tidak diperlukan}}}
\isian{Tempelkan dokumen pada halaman ini}

\judulhalaman{LAMPIRAN 12. SERTIFIKAT DAN DOKUMENTASI JALUR KOMPETISI}
\petunjuk{Hanya bila berlaku: sertifikat juara 1 sampai 3 BELMAWA, dokumentasi, dan surat keterangan (Pedoman TA butir 3.3.6). \tahap{\tidakperlu{Tidak diperlukan}}{\wajib{WAJIB bila berlaku}}}
\isian{Tempelkan dokumen pada halaman ini}

\judulhalaman{LAMPIRAN 13. PEMBAGIAN BAGIAN INDIVIDU DAN GILIRAN PRESENTASI}
\petunjuk{Isi untuk setiap anggota: judul Laporan TA individu, bagian proyek yang menjadi tanggung jawabnya, dan bagian presentasi serta demo yang dibawakan. Harus sama dengan Surat Pembagian Kerja Tim. \tahap{\opsional{Opsional}}{\opsional{Opsional (disarankan)}}}
\begingroup\footnotesize\singlespacing\renewcommand{\arraystretch}{1.6}
\begin{center}
\begin{tabular}{|>{\raggedright\arraybackslash}p{0.26\textwidth}|p{0.13\textwidth}|p{0.13\textwidth}|p{0.13\textwidth}|p{0.13\textwidth}|}
	\hline
	\textbf{Bagian} & \textbf{Ang. 1} & \textbf{Ang. 2} & \textbf{Ang. 3} & \textbf{Ang. 4}\\ \hline
	Judul Laporan TA individu & & & & \\ \hline
	Bagian proyek / modul & & & & \\ \hline
	Bagian presentasi & & & & \\ \hline
	Demo produk & & & & \\ \hline
\end{tabular}
\end{center}\endgroup

\judulhalaman{LAMPIRAN 14. FORM KENDALI BIMBINGAN}
\petunjuk{Minimal 4 pertemuan per semester, ditandatangani pembimbing atau Kartu Kendali dari TASYS (Pedoman TA Tabel 5.4 dan 5.5). \tahap{\wajib{WAJIB}}{\wajib{WAJIB}}}
\isian{Tempelkan dokumen pada halaman ini}

\judulhalaman{LAMPIRAN 15. HASIL UJI SIMILARITY}
\petunjuk{Laporan Turnitin lengkap atas laporan individu ini, \textit{similarity} $\leq$25\% dan kutipan langsung $\leq$10\% (Pedoman TA 5.7). \tahap{\tidakperlu{Tidak diperlukan}}{\wajib{WAJIB}}}
\isian{Tempelkan dokumen pada halaman ini}

\judulhalaman{LAMPIRAN 16. LOG SESI AI DAN RIWAYAT VERSI}
\petunjuk{Log sesi AI bila memakai AI, serta bukti riwayat versi laporan dan kode (Pedoman TA 4.2.1). \tahap{\tidakperlu{Tidak diperlukan}}{\wajib{WAJIB}}}
\isian{Tempelkan dokumen pada halaman ini}
